\section{Preliminaries: carries, the governor set, margin variants}\label{sec:prelim}

\begin{definition}[carry counters]\label{def:kappa}
For a prime $p$ and $m\ge1$ set
\[
\ka_p(m)\;:=\;\#\Big\{j\ge1:\ \fr{m/p^j}\ge\tfrac12\Big\}.
\]
The count is finite: $\fr{m/p^j}=m/p^j<\tfrac12$ once $p^j>2m$.
\end{definition}

\begin{proposition}[Kummer, via Legendre]\label{prop:kummer}
\statusP\quad
For every prime $p$ and $m\ge1$,
\[
v_p\binom{2m}{m}\;=\;\ka_p(m).
\]
\end{proposition}

\begin{proof}
By Legendre's formula,
$v_p\binom{2m}{m}=\sum_{j\ge1}\big(\fl{2m/p^j}-2\fl{m/p^j}\big)$, and for
any real $y$ one has $\fl{2y}-2\fl{y}=\mathbf 1[\fr{y}\ge\tfrac12]$.
\end{proof}

\begin{definition}[governor set and margin variants]\label{def:governor}
\[
\D\;:=\;\Big\{m:\ m\mid\binom{2m}{m}\Big\}
\;=\;\Big\{m:\ \ka_p(m)\ge e\ \text{for every prime power }p^e\,\|\,m\Big\},
\]
the equality by Proposition~\ref{prop:kummer} applied prime by prime.
For $\sigma\in\{+,0,-\}$ define the margin-modified carry counters
\[
\ka^{(\sigma)}_p(m)\;:=\;\#\Big\{j\ge1:\ \fr{m/p^j}\in I^{(\sigma)}_j\Big\},
\qquad
I^{(+)}_j=\big[\tfrac12,\,1-p^{-j}\big),\quad
I^{(0)}_j=\big[\tfrac12,\,1\big),\quad
I^{(-)}_j=\big[\tfrac12-p^{-j},\,1\big),
\]
and the margin governor sets
\[
D_0^{(\sigma)}\;:=\;\Big\{m:\ \ka^{(\sigma)}_p(m)\ge e
\ \text{for every prime power }p^e\,\|\,m\Big\},
\]
so $D_0^{(0)}=\D$ and $\Dp\subseteq\D\subseteq\Dm$.
\end{definition}

\begin{remark}
The margin sets differ from $\D$ only when some digit datum
$\fr{m/p^j}$ falls within $p^{-j}$ of the interval endpoints
$\tfrac12$ or $1$ at a prime power relevant to $m$'s governor condition.
Numerically the sets are nearly identical: on $m\le6\times10^4$ one has
$|\Dp|=|\D|$ exactly, and the density of $\Dm\setminus\D$ measured at
$10^9$ is $\approx 6\times10^{-5}$ (Appendix~\ref{app:numerics}).
Lemma~\ref{lem:margins} below asserts $\dens(D_0^{(\pm)})=c_1$.
\end{remark}

\begin{definition}[membership criterion for $\W$]\label{def:W1-criterion}
Since $\gcd(n,n-1)=1$, every prime divides at most one of $n$, $n-1$, and
\[
n\in\W
\iff
\begin{cases}
\ka_p(n)\ge e & \text{for every }p^e\,\|\,n,\\[2pt]
\ka_q(n)\ge e & \text{for every }q^e\,\|\,n-1.
\end{cases}
\]
Note carefully: both conditions constrain the digits \emph{of $n$} --- the
second at the primes of $n-1$. The whole difficulty of the problem is that
the second family of conditions couples the multiplicative structure of
$n-1$ to the digit expansion of $n$.
\end{definition}

\begin{lemma}[shift identity]\label{lem:shift}
\statusP\quad
Let $q^j$ be any prime power and $n\ge2$. If $q\mid n-1$, then
\[
\fr{\frac{n}{q^j}}
=\begin{cases}
\fr{\dfrac{n-1}{q^j}}+q^{-j}, &
\fr{\dfrac{n-1}{q^j}}<1-q^{-j}\quad(\text{no wrap}),\\[10pt]
\fr{\dfrac{n-1}{q^j}}+q^{-j}-1, &
\fr{\dfrac{n-1}{q^j}}\ge1-q^{-j}\quad(\text{wrap}).
\end{cases}
\]
In the wrap case $\fr{n/q^j}<q^{-j}$, so $\fr{n/q^j}\ge\tfrac12$ is possible
only when $q^j=2$, in which case $\fr{n/q^j}=0<\tfrac12$ anyway: \emph{a
wrapped position never produces a carry.}
\end{lemma}

\begin{proof}
$n/q^j=(n-1)/q^j+q^{-j}$ exactly; reduce mod $1$. (The hypothesis
$q\mid n-1$ is not needed for the identity, only for its later use; we
state it to fix the context.) For the final sentence: in the wrap case
$\fr{n/q^j}<q^{-j}\le\tfrac12$, with equality $q^{-j}=\tfrac12$ only for
$q^j=2$, and then $\fr{(n-1)/2}\ge1-\tfrac12$ forces $n-1$ odd, i.e.\
$\fr{n/2}=0$.
\end{proof}
